\chapter{Scope and method}

\section{What this report is about}
A multi-core quantum computer connects small processor cores over a network-on-chip. When a program moves a logical qubit
from one core to another, the move is carried out by teleportation, which consumes an entangled pair (an EPR pair)
shared between the two cores. Pre-generating those pairs before the program needs them can hide their generation latency,
but every core has only a few \emph{communication memories} in which pairs can wait, and many moves may need the same
core at once. This report studies that contention (\emph{congestion}) on a $4\times4$ mesh of cores.

The simulator used is \emph{ACE}, an adaptive-continuous-entanglement simulator built on SeQUeNCe v0.8.1. The first part of the
report documents what was found wrong with how ACE released requests to the network, how it was corrected, and how the
correction was verified. The second part reports the experiments that the corrected simulator makes possible:
latency, fidelity and waiting of ten benchmarks across communication-memory budgets (Experiment~1), the split of each
core's congestion into its own endpoint demand and relay traffic (Experiment~2), a study of ACE's memory cut-off
mechanisms, a study of whether offline routing can reduce congestion, and the same questions on a torus.

\section{Ground rules followed}
The following rules were fixed at the start and held throughout.
\begin{itemize}
  \item Milliseconds measured in ACE are never compared with milliseconds from the native SeQUeNCe backend.
  \item Results from the legacy and the corrected scheduler are never merged in one table.
  \item A fallback to on-demand generation is not a failure, but a dropped request is.
  \item Nothing is claimed as measured unless it was run; mechanisms that were inferred are labelled as such.
  \item Congestion-aware routing is not claimed to exist: no routing policy is simulated, only bounded structurally
        (Chapter~\ref{ch:routing}).
  \item Structural (trace-based) results are not called physical results.
  \item Physical parameters were not changed silently; the choice of parameters is an open decision
        (Appendix~\ref{app:params}).
\end{itemize}

\section{How the work was organised}
The work followed a gated plan, each gate committed separately on the branch \code{codex/ace-scheduler-correctness}:
\begin{description}
  \item[Gate 0] a bit-for-bit baseline of the legacy simulator (164 result rows) that must still reproduce after every change;
  \item[Gate 1] an execution event log, an invariant checker, microcases and an audit of the legacy simulator;
  \item[Gate 2] making two requests that look alike distinct bookings (reservation identity);
  \item[Gate 3] replacing per-core serialization by a per-qubit request manager with event-driven admission;
  \item[Gate 4] verification on all ten benchmarks, three strategies, three mapping seeds and several memory budgets;
  \item[Gate 5] tracing every physical parameter to its source;
  \item[Gate 6] rerunning the earlier 378-run budget study under the corrected scheduler.
\end{description}
